\documentclass{article}
\usepackage[utf8]{inputenc}
\usepackage{amsmath,amssymb}
\usepackage[most]{tcolorbox}
\usepackage{geometry}
\geometry{margin=2.5cm}

\newcommand{\step}[1]{\par\noindent\hangindent=3.5em\hangafter=1%
  \makebox[3.5em][l]{\textbf{#1.}}}
\newcommand{\steptag}[1]{\hfill{\small\textsf{[#1]}}}

\newtcolorbox{proofbox}[1]{
  colback=gray!5, colframe=gray!60,
  fonttitle=\bfseries, title={#1},
  breakable, enhanced,
  left=4pt, right=4pt, top=4pt, bottom=4pt,
  before skip=10pt, after skip=10pt
}

\begin{document}

\begin{proofbox}{Mass split: for an exact signed idempotent P and any row ...}
\step{1} Mass split: for an exact signed idempotent P and any row index v, writing a\_j = P\_\{vj\}, a\_j\textasciicircum{}+ = max(a\_j, 0), a\_j\textasciicircum{}- = max(-a\_j, 0), and nu\_v = sum\_j a\_j\textasciicircum{}-, one has sum\_j a\_j\textasciicircum{}+ = 1 + nu\_v. \steptag{validated} \\[6pt]
\step{1.1} By def-signed-idempotent, P 1 = 1; therefore for the chosen row index v, sum\_j P\_\{vj\} = 1, and since a\_j = P\_\{vj\}, sum\_j a\_j = 1. \steptag{validated} \\[4pt]
\step{1.2} For each index j, the definitions a\_j\textasciicircum{}+ = max(a\_j,0) and a\_j\textasciicircum{}- = max(-a\_j,0) imply a\_j = a\_j\textasciicircum{}+ - a\_j\textasciicircum{}-. \steptag{validated} \\[6pt]
\step{1.2.1} For a fixed j with a\_j >= 0, max(a\_j,0)=a\_j and max(-a\_j,0)=0, so a\_j\textasciicircum{}+ - a\_j\textasciicircum{}- = a\_j - 0 = a\_j. \steptag{validated} \\[4pt]
\step{1.2.2} For a fixed j with a\_j < 0, max(a\_j,0)=0 and max(-a\_j,0)=-a\_j, so a\_j\textasciicircum{}+ - a\_j\textasciicircum{}- = 0 - (-a\_j) = a\_j. \steptag{validated} \\[4pt]
\step{1.2.3} Every real number a\_j satisfies exactly one of a\_j >= 0 or a\_j < 0; hence the identity a\_j = a\_j\textasciicircum{}+ - a\_j\textasciicircum{}- holds for every index j. \steptag{validated} \\[4pt]
\step{1.3} Summing the pointwise identity from child 1.2 over the finite index set gives sum\_j a\_j = sum\_j a\_j\textasciicircum{}+ - sum\_j a\_j\textasciicircum{}-. \steptag{validated} \\[4pt]
\step{1.4} By the definition of nu\_v in the root statement, nu\_v = sum\_j a\_j\textasciicircum{}-. \steptag{validated} \\[4pt]
\step{1.5} Combining child 1.1 with child 1.3 and child 1.4 gives 1 = sum\_j a\_j\textasciicircum{}+ - nu\_v; adding nu\_v to both sides gives sum\_j a\_j\textasciicircum{}+ = 1 + nu\_v. \steptag{validated} \\[4pt]
\end{proofbox}

\end{document}
